\section{Conclusion and Future Work}
\label{dyn_sec:conclusion}

We presented \ourmethod, a method for articulated dynamic object reconstruction from sparse temporal observations. \ourmethod{} leverages a rough input skeleton and an initial static reconstruction to learn a skeleton-driven deformation field that models coherent motion across time.
Furthermore, we showed that the need for multi-view initialization can be relaxed using a pre-trained diffusion-based generative prior, enabling dynamic reconstruction in real-world scenarios.
Experiments on synthetic datasets show that \ourmethod{} outperforms existing methods by up to 34\% in PSNR under sparse observations and performs comparably to dense monocular methods on real-world datasets, even though \ourmethod{} uses $10\times$ fewer frames. %\vtxt{clarity this sentence. who has far fewer frames?}
While promising, our approach has limitations. The diffusion-based initialization can fail under severe self-occlusion or uncommon viewpoints, as it relies on a general pre-trained model. 
Moreover, test-time interpolation may struggle with highly complex motion. 
A potential future direction is to investigate using  category-specific priors or a pre-trained prior conditioned on the noisy skeleton input to guide motion estimation and reconstruction.
 
